% !TeX root = ../navier-stokes-manuscript.tex
% ============================================================
% 2.5 Simultaneous Sobolev Intersections
% ============================================================

\subsection{Simultaneous Sobolev Intersections}\label{sub:SimultaneousSobolevIntersections}

The critical-height calculation reaches spatial energies of arbitrarily high order, but any one derivative request stops at a finite order.
Suppose, for example, that one wants \(\partial_x^\alpha u\) to be continuous for every \(\lvert\alpha\rvert\le k\).
In three dimensions, the Sobolev embedding
\[
  H^m(\R^3)
  \hookrightarrow
  C^k(\R^3)
\]
holds whenever
\[
  m>k+\frac32.
\]
One requested spatial order \(k\) comes from one finite height \(m\) above its embedding threshold.
As \(k\) ranges through all finite orders, no last \(m\) can be chosen, but no infinite-order norm is required either.
The exact spatial object is the coordinatewise intersection
\[
  \bigcap_{m=0}^{\infty}H^m(\R^3),
\]
whose membership means that every finite Sobolev norm is finite on the same field.
For each fixed \(k\), one of those finite memberships yields the desired \(C^k\) control, so
\[
  \bigcap_{m=0}^{\infty}H^m(\R^3)
  \hookrightarrow
  C^\infty(\R^3).
\]

Finite-time response asks for the same construction in the temporal direction.
Let \(I=(0,T)\) with \(T<T_*\), and let \(X\) be a Hilbert space.
The vector-valued embedding
\[
  H^r(I;X)
  \hookrightarrow
  C^q(\overline I;X)
\]
holds whenever
\[
  r>q+\frac12.
\]
One requested time derivative of order \(q\) comes from one finite temporal height \(r\).
Allowing every finite request gives
\[
  \bigcap_{r=0}^{\infty}H^r(I;X)
  \hookrightarrow
  C^\infty(\overline I;X).
\]
The time intersection is again coordinatewise: each \(q\) chooses its own finite \(r\).

The mixed jet requires these choices at the same time.
Ask for one concrete coordinate,
\[
  \partial_t^2\partial_x^\alpha u,
  \qquad
  \lvert\alpha\rvert\le4.
\]
The single finite membership
\[
  u\in H^3\bigl(I;H^6(\R^3)\bigr)
\]
places that coordinate in
\[
  C^2\bigl(\overline I;C^4(\R^3)\bigr),
\]
because
\[
  3>2+\frac12,
  \qquad
  6>4+\frac32.
\]
Nothing special happened at the pair \((2,4)\).
For arbitrary finite temporal and spatial orders \(q\) and \(k\), choose finite indices satisfying
\[
  r>q+\frac12,
  \qquad
  m>k+\frac32.
\]
Then
\[
  H^r\bigl(I;H^m(\R^3)\bigr)
  \hookrightarrow
  C^q\bigl(\overline I;C^k(\R^3)\bigr).
\]

To satisfy every finite time--space request on the same velocity, take the intersection of these finite Sobolev memberships:
\[
  \mathscr V_\infty(I)
  :=
  \bigcap_{r=0}^{\infty}
  \bigcap_{m=0}^{\infty}
  H^r\bigl(I;H^m(\R^3)\bigr).
\]
Membership means
\[
  u\in\mathscr V_\infty(I)
  \quad\Longleftrightarrow\quad
  \norm{u}_{H^r(I;H^m)}
  <
  \infty
  \quad
  \text{for every finite pair }(r,m).
\]
Each membership is a coordinate of the same \(u\), and the intersection keeps their norms separate.
The embeddings above then give
\[
  \mathscr V_\infty(I)
  \hookrightarrow
  C^\infty\bigl(\overline I\times\R^3\bigr).
\]

Choose temporal and spatial indices large enough to give a continuous trace above the critical levels.
That trace contains both
\[
  u(t)\in H^{1/2}(\R^3)
  \qquad\text{and}\qquad
  u(t)\in H^{3/2}(\R^3)
\]
at each time in the closed interval.
The critical quantities are then defined on the closed interval by
\[
  H(t)=\frac12\norm{\Lambda^{1/2}u(t)}_2^2,
  \qquad
  E_{3/2}(t)=\frac12\norm{\Lambda^{3/2}u(t)}_2^2
\]
The calculation of critical height led to these spatial levels; the intersection is what controls them at the endpoint.

Pressure is not an element of the velocity intersection.
Write
\[
  Q_n=\partial_t^np,
\]
at every finite depth.
The pressure Poisson equation recovers each \(Q_n\) from the velocity rows, while the differentiated momentum equation keeps \(\nabla Q_n\) coupled to \(V_{n+1}\), viscosity, and transport.
Thus the pressure jet is recovered from \(u\) without treating pressure as an element of the velocity space.

Each estimate controls one finite block of this intersection.
The two Sobolev levels needed to produce a common time trace determine the upper and lower edges of that block.
